\chapter*{Conferencia III}
\addcontentsline{toc}{chapter}{LECTURA I - El timbre y el carácter del violín}
CABALLEROS: Ustedes comprenden el objetivo de nuestra
La sociedad incluye el estudio tanto del sonido musical como del no musical producido por el violín. No es
Resulta aceptable admitir que el violín pueda tener un sonido distinto al musical. Instintivamente, solemos considerar "dulzura" y "violín" como sinónimos. Incluso hablar de un violín "ruidoso" resulta irritante para nuestros sentidos, mientras que escucharlo a la fuerza es una auténtica tortura. Bajo el epígrafe de "Crueldad contra la Humanidad", la venta y el uso de violines "ruidosos" deberían estar prohibidos por ley. Es lamentable que el sonido más perjudicial para la música provenga de imitaciones del violín. La apariencia de estas imitaciones fraudulentas es la trampa tendida al comprador inocente. Debido a las apariencias, el comprador cree estar en posesión de un violín. Hoy en día, encontramos víctimas de este engaño por doquier. En los oídos crédulos de cada una de estas víctimas resonaba aquel canto de sirena: "El violín siempre mejora con la edad y el uso". Al intentar aplicar esta vieja canción a las imitaciones del violín, comienzan los problemas. La familia de la víctima sufre un ataque de nervios y su entorno inmediato se ve afectado.
es completamente evitada por la música. Su laborioso "cobrar" se ve impulsado a una mayor actividad por la esperanza de mejorar su tono. Con firmeza y afecto, se dirige a esa imitación como "vieja concha". Los años pasan sin ninguna mejora en el tono. Con frecuencia, personas mayores, con dudas, me han traído estas imitaciones. (Con dudas), "Doctor, ¿puede...?"
¿Me dices qué le pasa a mi vieja carcasa? Rápidamente paso el "sonido" doblado por la superficie interior de la caja de resonancia, el traqueteo que produce me recuerda al carro del hojalatero en el camino de pana.
—Sí, señor, puedo decírselo. —¿Qué es? —Es un fraude.
Al retirar la caja de resonancia, permitiendo que sus ojos, llenos de dudas, contemplen la escena dentro de su "vieja coraza", la sonrisa congelada en su rostro arrugado resulta verdaderamente patética. La graduación (disculpen la expresión) de esta caja de resonancia de imitación se realiza completamente con gubia y navaja. La barra es un "trozo" que queda tras la prisa de la gubia, y cuidadosamente alisado con la navaja en el lado expuesto a su inspección. Como no puede verlos, los dos superiores
Los rincones quedan olvidados.	Los revestimientos, de varios di
Las menciones son a la vez demasiado largas y demasiado cortas.	Afuera
A simple vista, este fraude se parece al violín. Aunque el barniz es de mala calidad, el color es bastante...
artístico.
Aquí hay toda una vida dedicada a la música, pero desperdiciada, desperdiciada por esta maldita imposición a un comprador inocente, crédulo e incauto. Además, en este gran e ilustrado (perdón una vez más, )
Estados Unidos de América del Norte, que otorga libertad, vende miles de copias de este mismo fraude premeditado y descarado. Semejante imposición a los amantes de la música es motivo suficiente para despertar sentimientos homicidas en un corazón de piedra. Estos fraudes son producto de la codicia, la estupidez y la maldad combinadas.
falta. La codicia, en estos fraudes, se demuestra al cobrar más que el precio de la leña, a pesar de que las cotizaciones se colocan en la columna de "pfennige". En estos fraudes, la estupidez se demuestra al pensar que quince minutos más de trabajo en la caja de resonancia podrían colocar las cotizaciones en la columna de "marks". En estos fraudes, la crueldad se demuestra con el asesinato premeditado de Music. La criminalidad de estos estafadores supera la de quien roba a niños pequeños, porque combinan el asesinato con el robo. ¡Pobre Music! ¡Horrorizada al ver su lugar favorito ya no habitable! Mientras profesas amistad, tú, Sr. Estafador, eres la causa de su angustia. A ti, y a los como tú, y de ti, digo: "El infierno no conoce a un villano mayor".
¡Tú, violín! ¡Tú!
¡Quién es el dulce Rey de la Música!
¿Debe tu forma hechicera ocultar el punzante y mortal aguijón de la codicia? ¡Hombres! ¡Levántense! ¡Pónganse de pie ahora en filas de aquellos que desean luchar para matar a esta criatura ladrona y mortal, y darle a la Música lo que le corresponde por derecho!
"Doctor, ¿cómo puede un comprador inocente reconocer la mano de la codicia?"
Amigo mío, el caso es desesperado cuando el comprador es tan inocente que, por ello, es incapaz de leer "Hecho en Alemania".
Después de toda una vida dedicada a la producción de violines-
En cuanto al tono, me veo obligado a reconocer la incredulidad ante la posibilidad de una solución satisfactoria para todos los problemas que conlleva. También me veo obligado a afirmar que mi experiencia y conclusiones sobre cómo funciona el violín para producir sonido musical no corroboran las teorías de Savart, Honeyman y otros. Si bien no puedo resolver todos los problemas relacionados con el tono del violín, no considero que mis esfuerzos en este sentido sean completamente infructuosos. Les presento mis conclusiones y principios, en la medida en que me resultan satisfactorios, en términos tan claros y concisos como me es posible. No me detendré en teorías abstractas. Todas mis conclusiones se basan en pruebas prácticas y repetidas realizadas con violines usados, excepto los principios de graduación de la caja de resonancia para lograr la máxima uniformidad del tono. Como se demostrará más adelante, este principio, debido a una discapacidad física, solo recibió una demostración. Les lego a ustedes todo aquello de valor que puedan contener mis conclusiones.
Ahora me referiré al tema del tono del violín. [Esta característica del tono del violín merece una cuidadosa consideración. Siendo satisfactorias otras cosas, el tono debe determinar el lugar de utilidad de cada violín. Así pues: El violín de tono grave en toda la orquesta no debe colocarse al frente de la orquesta completa, y, debido a que sus tonos no tienen suficiente intensidad (potencia de proyección) para satisfacer al oído lejano. El lugar del violín de tono grave en toda la orquesta es el de instrumento solista. Asimismo, el violín de tono agudo
En general, no debería emplearse como instrumento solista, ya que sus tonos agudos, por sí solos, resultan desagradables al oído. He visto cómo la reputación de un solista de renombre se arruinaba por utilizar un buen violín de orquesta, de tono agudo, como instrumento solista.
He formulado ocho reglas, cada una de las cuales permite modificar la afinación del violín. Para facilitar su consulta, estas reglas están numeradas consecutivamente del I al VIII. Dado que se hará referencia a ellas con frecuencia, se agrupan.
Regla /.Alargar un agente productor de tono, manteniendo las demás dimensiones iguales, disminuye el tono.
Regla //.Acortar un agente productor de tono, manteniendo las demás dimensiones iguales, eleva el tono.
Regla III. Al aumentar el grosor de un agente productor de tono, manteniendo las demás dimensiones iguales, aumenta el tono.
Regla IV. Al disminuir el grosor de un agente productor de tono, manteniendo las demás dimensiones iguales, disminuye el tono.
Regla V. El alargamiento de las columnas de aire perpendiculares confinadas reduce el tono.
Regla VI. Acortar las columnas de aire confinadas y perpendiculares eleva el tono. 
Regla VII. Ampliar las salidas de la caja de resonancia eleva el tono. 
Regla VIII. Reducir las salidas de la caja de resonancia disminuye el tono.
En la construcción práctica de violines, o en la afinación de violines,
\label{plancha_fondo}
Cada una de estas reglas tiene su valor. En mi práctica, las aplico únicamente a la caja de resonancia. De ninguna manera las aplico a la tapa trasera. Tras una larga y continua experimentación, he llegado a la conclusión de que la función principal de la tapa trasera es la de un medio reflectante. Por lo tanto, solo exijo de la tapa trasera la rigidez suficiente para soportar, sin temblor, la tremenda carga de movimiento de ondas moleculares que se origina en la superficie interna de la caja de resonancia; y que su fibra sea fina, densa y susceptible de un pulido intenso. Así, una vez terminada, la superficie interna de la tapa trasera presenta una superficie reflectante perfecta. No digo que la tapa trasera no pueda modificar el tono. Al contrario, afirmo que la rigidez de la tapa trasera puede reducirse hasta que el tono del violín se vuelva hueco, débil e inservible. He encontrado muchos violines arruinados de esta manera; más aún que por una gran reducción de la rigidez de la caja de resonancia. De hecho, muchos constructores de violines intentan regular el tono reduciendo la rigidez de la tapa trasera. En mi experiencia con la prueba de larga distancia para la intensidad (poder de transmisión) del tono, esta manera de regular el tono ha demostrado ser seriamente errónea. El tono de cada violín, regulado de esta manera, ha quedado muy por debajo del récord de distancia. En mi opinión, la regulación del tono mediante la reducción de la rigidez de la caja de resonancia ha demostrado ser el método más seguro. Por lo tanto, estas reglas para controlar el tono se aplican a la
caja de resonancia.
Ahora consideraré la aplicación práctica de la Regla I: "Alargar un agente productor de tono, manteniendo las demás dimensiones iguales, disminuye el tono". En la aplicación de esta "regla" a la caja de resonancia, sus efectos son mayores sobre el tono de las cuerdas G y D que sobre el tono de las cuerdas A y E. Como ayuda para la ilustración, supondré que aquí hay una caja de resonancia con un grosor uniforme de 8,64. Este grosor de la caja de resonancia, junto con las dimensiones de las barras, en longitud,
lOi pulgadas, grosor, 3-16, profundidad, I, que se estrecha hacia los extremos, dan un carácter agudo al tono de las cuerdas G y D, no el más agudo, pero sí agudo. Deseando bajar el tono de las cuerdas G y D, procedo a alargar el agente productor de tono según la Regla I. Con el propósito de limitar el tono a las cuerdas G y D, a media pulgada a cada lado de la barra dibujo líneas a lápiz que se extienden lo más cerca posible del fileteado y los bloques de los extremos para reducir el grosor de la caja de resonancia. Dentro de estas líneas, comenzando en la posición del puente, reduzco gradualmente el grosor hasta 4-64 en los extremos de la placa; también, fuera de estas líneas, el grosor aumenta gradualmente hasta 8-64. Al probar, se encuentra que el tono de las cuerdas G y D se ha bajado perceptiblemente. De este modo se demuestra que, en la caja de resonancia de 8-64 en toda su longitud, no actúa con la energía suficiente para producir un sonido audible; asimismo, de este modo se demuestra la corrección de la Regla I.
Llamo a la Regla II: "Acortar un tono que produce
agente, manteniendo las demás dimensiones iguales, plantea
tono-tono. ' ' Evidentemente, la aplicación de esta regla debe ocurrir en el momento de construir la caja de resonancia; después, no puede aplicarse. En la ilustración para bajar el tono, al alargar el agente productor de tono, hay datos suficientes para subir el tono al acortar el agente productor de tono.
agente.
Regla III: "Aumentar el grosor de un agente productor de tono, manteniendo las demás dimensiones iguales, eleva el tono." Esta regla se aplica en el momento de la construcción. He encontrado algunas cajas de resonancia con un grosor de 1 pulgada en la posición del puente; y dos con el mismo gran grosor a lo largo de la unión central. El tono extremadamente agudo de la caja de resonancia tan pesada en madera no se disfruta ni siquiera a distancia al aire libre. En la época de los trovadores, este tono agudo al aire libre era popular; de hecho, bastante necesario para satisfacer el gusto del público. Hoy en día, el músico callejero, imitando "los viejos tiempos", todavía tiene éxito en hacer que el público se deshaga de sus monedas mediante el uso de violines de tono tan agudo; y personalmente, no me opongo, siempre que acepte mi moneda como una señal para que se vaya. (Honeyman, de nuevo, es autoridad para la afirmación de que la caja de resonancia de Nicolas Amatus tiene un grosor de 1 pulgada a lo largo de la unión central, y de ahí, lateralmente, hasta en el bordes. )
Regla IV: "Disminuir el grosor de un agente productor de tono, manteniendo las demás dimensiones iguales, reduce el tono." El tono agudo de la caja de resonancia gruesa se reduce al disminuir
su grosor. Disminuir el tono reduciendo el grosor de la caja de resonancia requiere precaución, porque este trabajo pone en práctica dos factores poderosos, como se mostrará después de la Regla V, a saber:
"Alargar las columnas de aire perpendiculares confinadas reduce el tono." Así, al reducir el tono disminuyendo el grosor de la caja de resonancia, esta misma reducción también alarga las columnas de aire perpendiculares dentro del cuerpo del violín; por lo tanto, intervienen dos factores; y de ahí la necesidad de precaución. En una caja de resonancia de madera resonante, con el grosor reducido a 8,64 en toda su extensión, una reducción adicional de 1,64 provoca una disminución perceptible del tono; y, a partir de 7,64, la reducción adicional de 1,64 reduce el tono dos veces más que la disminución producida al pasar de 8,64 a 7,64. Atribuyo este aumento aritmético en la disminución del tono a la acción simultánea de dos factores.
Regla VI: "Acortar las columnas de aire perpendiculares y confinadas eleva el tono". Para elevar el tono de un violín cuyas tapas tienen un grosor igual o inferior a los valores de peligro, he aplicado con éxito esta regla disminuyendo la profundidad de las costillas; acortando así las columnas de aire perpendiculares dentro del cuerpo. El hecho de que acortar las columnas perpendiculares eleve el tono se demuestra de numerosas maneras conocidas, como en el tubo del órgano, el silbato de vapor, etc. Como demostración práctica de esta regla, he disminuido la profundidad de las costillas en 1/16 de pulgada, realizando una prueba en cada reducción. El resultado proporciona una prueba fehaciente de la veracidad de la Regla VI.
Regla VII: «Aumentar el tamaño de las salidas eleva el tono». En la práctica, solo aplico esta regla después de probar el tono de cada violín. La experiencia confirma su veracidad. Por seguridad, las áreas de las salidas deben ser relativamente pequeñas al construir la caja de resonancia. Si, tras probar el tono del violín terminado, este resulta ligeramente inferior al deseado, se puede corregir aumentando ligeramente el área de las salidas. En este trabajo, prefiero realizar dicho aumento retirando madera del borde interior de las salidas, con el fin de acercarlas a los puntos focales de concentración de ondas sonoras. [Más adelante, se analizará la concentración de ondas sonoras en las salidas.]
Regla VIII: "Reducir el área de salidas disminuye el tono." Esta regla debe aplicarse durante la construcción, aunque también es posible aplicarla al violín terminado. Reducir el área de salidas es un factor importante para disminuir el tono y el volumen. Para obtener el sonido de violín más agradable para uso en estudio, no conozco ningún modificador del tono de mayor importancia que la pequeña
salida.
El sonido musical es un tema fascinante para el estudiante de violín. En esa cualidad del sonido llamada 'tono'
Existe abundante material para el estudio. Algunos problemas relacionados con el sonido tienen solución; otros no. El problema del tono absoluto se puede resolver mediante mecanismos capaces de registrar el número de golpes.
por segundo transmitido a las moléculas de aire. Por lo tanto, el tono de cualquier sonido está determinado por el agente que lo produce. Sin embargo, este hecho no resuelve todos los problemas relacionados con el tono. El tono absoluto no lo abarca todo. [*]Las palabras parecen insuficientes para siquiera enumerar todos los problemas en el tono del sonido.. Al observar los problemas de afinación desde el punto de vista del violín, nos vemos obligados a admitir ocho factores modificadores (tal como se formulan en las Reglas I-VIII) en la producción del sonido del violín, mientras que solo uno de estos factores es el
¡Agente impactante!
Formular los problemas del sonido del violín es un trabajo arduo; resolverlos después de formularlos es aún más difícil. Les dejo la solución. Por sus miradas de reojo y sus rostros esquivos, deduzco que no consideran mi legado un bien valioso. Quizás tal bien no lo sea. Es fácil legar bienes de los que uno no es dueño. Pero, dado que la solución al problema del sonido del violín no se encuentra ahora, no es seguro que nunca se encuentre. Les aconsejo que continúen esforzándose. Poco a poco, la solución a estos problemas llegará a cada estudiante perseverante.
Para continuar con el tema, debo recurrir a supuestos y combinarlos con algunos hechos concretos, junto con razonamientos por analogía. Que esta combinación les resulte agradable o no dependerá en gran medida de su entusiasmo y de mi capacidad de persuasión.
Procedencia: Un cuerpo puede vibrar y, sin embargo, no \label{sonido_audible}
producir sonido audible. Para producir sonido audible, el agente productor de sonido debe aplicar golpes de energía suficiente sobre moléculas de aire contiguas, de manera que provoquen movimiento en cada molécula de aire entre el punto de impacto y nuestros tímpanos; de lo contrario, no percibimos el sonido. Estos son hechos. El siguiente punto es una suposición. Por lo tanto: supongo que la longitud total de 14 pulgadas de la caja de resonancia ensamblada no puede actuar sobre las moléculas de aire con la energía suficiente para producir sonido audible; pero me veo obligado a conformarme con una aproximación. Las dificultades para lograr precisión en este asunto son grandes. Las diferencias en la elasticidad de diferentes muestras de madera son infinitas. Además, un ligero aumento en la altura del arco incrementa la rigidez de la caja de resonancia. Como mera suposición, considero que la longitud máxima de la caja de resonancia que produce sonido audible es de 12 pulgadas. En la práctica, para asegurar dicha longitud de 12 pulgadas, reduzco el grosor desde la posición del puente gradualmente hasta 4,64 lo más cerca posible de los extremos de la caja de resonancia. Considero que el grosor de la caja de resonancia en los bordes, cuando se reduce a 4-64, está en el límite de la seguridad; y este grado de reducción solo lo empleo cuando deseo el tono más bajo y práctico para las cuerdas de sol.
Si ocurriera que una reducción excesiva del tono sigue a una reducción del grosor de la caja de resonancia, una reducción del grosor suficiente para causar un tono hueco y débil, no conozco ningún remedio que consista en aumentar el grosor de la caja de resonancia; es decir, ningún remedio de este tipo tiene valor. Sí sé que en
En la ciudad de Chicago vive un hombre que afirma tener la capacidad de restaurar el tono de los violines mediante una fina chapa de madera aplicada a la superficie interior de la caja de resonancia. Me trajeron un violín así tratado con la petición de que hiciera algo por su sonido apagado. Sin saber qué se le había hecho a este violín, primero apliqué el arco como de costumbre, esperando así determinar el trabajo necesario. Pero su sonido me quitó por completo la idea. De ninguna manera podía predeterminar la causa de tan extraordinaria falta de brillo. Con prisa
¡Eliminó la mesa de resonancia - Chicago!
¡Tú, tan cuidadoso de no ocultar tus talentos! ¡Cómo has perdido tu prestigio! Has cuidado bien de tus cirujanos. Su reputación es mundial. Pero la reputación de tu cirujano de violines no ha sido igualmente cuidada. Con tu permiso, ayudaré a corregir este descuido. Primero, describiré la intervención de cirugía plástica en este violín. Segundo, prescribiré un tratamiento para el cirujano. El área cubierta por la chapa equivale a la mitad de la caja de resonancia; sin embargo, estas áreas están divididas en dos territorios, uno en cada extremo de la placa. El grosor de la chapa es de 1,32. Su color es marrón oscuro. Su textura es similar a la del papel grueso. Su resistencia no es comparable a la del buen papel. Está cortada en piezas de dimensiones variadas; algunas miden 2x4 pulgadas, otras tan solo 2 pulgadas. En algunos lugares, la chapa está doblada. Al colocar los fragmentos, no se sigue ningún orden específico; la longitud de algunos se extiende a través de la veta.

de la caja de resonancia; de otras, en la dirección de la veta, y otras, oblicuas a la veta, mientras que entre parches hay áreas inexplicablemente pasadas por alto. La escena provoca un asombro sin límites. Hasta este momento, me enorgullecía de mi capacidad para determinar la causa de los defectos de sonido del violín mediante el uso del arco.
Ahora mi orgullo está humillado. Se dice que uno nunca...
Se vuelve demasiado viejo para aprender. Es cierto. ¡Pobre violín! De tu degradación estoy aprendiendo. Antes de conocer tu lamentable condición, supuse que esos monstruos insaciables, el calor y la humedad, eran tus enemigos más implacables. Pero has venido a enseñarme que estaba equivocado, a enseñarme que tu peor enemigo es el hombre. ¡Y este hombre vive en Chicago! Estos parches de chapa no están unidos a la caja de resonancia con pegamento, sino con una pasta espesa. ¡Gracias por la pequeña piedad! Con un paño escurrido de agua caliente, la pasta se ablanda fácilmente y levanto los parches enteros. Esos parches, donde la chapa está doblada, me llenan de sorpresa por la generosidad de un habitante de Chicago. Su factura podría ser igual de grande si se ahorrara las piezas sobrantes. Todo el asunto es inexplicable. Cómo este hombre puede vivir en Chicago y escapar de la notoriedad es uno de los misterios. Pero la receta, merecidamente justificada por este hombre, no tiene nada de misterioso. Es una receta dictada por la justicia. Aquí hay un violín sin defectos de grado grave, excepto el del modelo. El arqueado cae en los extremos de las tapas demasiado abruptamente para la potencia del sonido; por lo demás es un violín bastante bueno. Su caja de resonancia es realmente una selección superior. Después de haber sido nuevamente en condiciones de tocar, yo
Encuentran que este violín vale \$75. Pero, su valor, como provenía de los trabajos de chapado, era menos de setenta y cinco centavos. ¡Amo la música! ¡Dulce música! Me deleito en rendirle homenaje, incluso hasta el grado de adoración. Por lo tanto, cuando veo manos asesinas extendiéndose hacia la música, el león dentro de mí se despierta. Lleva a este miserable culpable al lago, (no tan lejos como para alcanzar agua pura), átale un peso al cuello, átale un salvavidas a los pies,
(de lo contrario, esa cabeza no se hundirá), tírenlo por la borda, manténganlo bajo el agua hasta que, por el cese de las burbujas que suben, sea seguro que dentro de su cráneo hay un incremento de materia gris.
